\documentclass[12pt,a4paper]{article}

\usepackage[utf8]{inputenc}
\usepackage{amsmath}
\usepackage{amssymb}
\usepackage{geometry}
\usepackage{booktabs}

\geometry{top=1in, bottom=1in, left=1in, right=1in}

\title{\textbf{Lecture 18: Marginal PMF, Correlation, Covariance, Joint Gaussian RVs}}
\author{}
\date{}

\begin{document}

\maketitle

\section{Marginal PMF}

\subsection{Definition}
For two discrete random variables $X$ and $Y$, the joint PMF is defined as:
\[
p_{XY}(x,y) = P(X=x, Y=y)
\]

\textbf{Key Idea (from lecture)} \\
The marginal PMF of one variable is obtained by summing the joint PMF over all possible values of the other variable.

\subsection{Step-by-Step Derivation}
To extract marginal distributions:

\textbf{For $X$:}
\[
p_X(x) = \sum_{y} p_{XY}(x,y)
\]
\textit{Reasoning:}
\begin{itemize}
    \item Fix $x$
    \item Sum probabilities over all possible $y$
    \item This accounts for all outcomes where $X=x$
\end{itemize}

\textbf{For $Y$:}
\[
p_Y(y) = \sum_{x} p_{XY}(x,y)
\]
\textit{Reasoning:}
\begin{itemize}
    \item Fix $y$
    \item Sum over all $x$
    \item This captures all events where $Y=y$
\end{itemize}

\subsection{Marginal PMF from Joint PMF Table}
From the table (page 4):
\begin{itemize}
    \item Row sums $\rightarrow p_X(x)$
    \item Column sums $\rightarrow p_Y(y)$
\end{itemize}

\textbf{Mathematically:}
\[
p_X(x_i) = \sum_j p_{ij}, \quad p_Y(y_j) = \sum_i p_{ij}
\]

\textbf{Normalization Property:}
\[
\sum_i p_X(x_i)=1, \quad \sum_j p_Y(y_j)=1
\]

\subsection{Worked Example (Step-by-Step)}
From the table (page 5):

\begin{table}[h!]
\centering
\renewcommand{\arraystretch}{1.5}
\begin{tabular}{|c|c|c|}
\hline
$X \setminus Y$ & $y=0$ & $y=1$ \\
\hline
$x=0$ & $1/8$ & $1/4$ \\
\hline
$x=1$ & $3/8$ & $1/4$ \\
\hline
\end{tabular}
\end{table}

\textbf{Step 1: Compute $p_X(x)$} \\
For $x=0$:
\[
p_X(0) = \frac{1}{8} + \frac{1}{4} = \frac{1}{8} + \frac{2}{8} = \frac{3}{8}
\]
For $x=1$:
\[
p_X(1) = \frac{3}{8} + \frac{1}{4} = \frac{3}{8} + \frac{2}{8} = \frac{5}{8}
\]

\textbf{Step 2: Compute $p_Y(y)$} \\
For $y=0$:
\[
p_Y(0) = \frac{1}{8} + \frac{3}{8} = \frac{4}{8} = \frac{1}{2}
\]
For $y=1$:
\[
p_Y(1) = \frac{1}{4} + \frac{1}{4} = \frac{1}{2}
\]

\textbf{Final Marginals}
\[
p_X(x) =
\begin{cases} 
3/8, & x=0 \\ 
5/8, & x=1 
\end{cases}
\quad
p_Y(y) =
\begin{cases} 
1/2, & y=0 \\ 
1/2, & y=1 
\end{cases}
\]

\section{Correlation}

\subsection{Definition}
Correlation is introduced via the product moment:
\[
R_{XY} = E[XY]
\]

\subsection{Interpretation (Step-by-Step)}
From lecture (pages 6–7):
\begin{itemize}
    \item Measures average value of product $XY$
    \item Depends on:
    \begin{itemize}
        \item Mean values
        \item Spread
        \item Dependence
    \end{itemize}
\end{itemize}

\textbf{Important Observation}
\[
R_{XY} = 0 \Rightarrow \text{orthogonal about origin}
\]

\subsection{Limitation}
\begin{itemize}
    \item Computed about the origin (not means)
    \item Mixes: mean, variance, dependence
\end{itemize}
$\Rightarrow$ Therefore, it is not a clean measure of dependence.

\textbf{Transition to Covariance:} To remove mean effects:
\[
\text{Center variables } \Rightarrow \text{Covariance}
\]

\section{Covariance}

\subsection{Definition}
\[
\text{Cov}(X,Y) = E[(X-\mu_X)(Y-\mu_Y)]
\]

\textbf{Step-by-Step Meaning}
\begin{itemize}
    \item \textbf{Subtract means:} $(X - \mu_X)$ and $(Y - \mu_Y)$
    \item \textbf{Multiply deviations:} Measures joint variation
    \item \textbf{Take expectation:} Average behavior
\end{itemize}

\subsection{Interpretation}
From page 9:
\begin{itemize}
    \item $\text{Cov}(X,Y) > 0$: same direction
    \item $\text{Cov}(X,Y) < 0$: opposite direction
    \item $\text{Cov}(X,Y) = 0$: no linear relation
\end{itemize}
\textbf{Warning:} Magnitude depends on units $\rightarrow$ cannot measure strength directly.

\subsection{Algebraic Derivation (Step-by-Step)}
\begin{align*}
\text{Cov}(X,Y) &= E[(X-\mu_X)(Y-\mu_Y)] \\
&= E[XY - X\mu_Y - Y\mu_X + \mu_X\mu_Y] \quad \text{(Expand product)} \\
&= E[XY] - \mu_Y E[X] - \mu_X E[Y] + \mu_X\mu_Y \quad \text{(Apply expectation linearly)} \\
&= E[XY] - \mu_X\mu_Y - \mu_X\mu_Y + \mu_X\mu_Y \quad \text{(Substitute } E[X]=\mu_X, E[Y]=\mu_Y \text{)} \\
&= E[XY] - \mu_X\mu_Y
\end{align*}

\textbf{Final Formula}
\[
\boxed{\text{Cov}(X,Y) = E[XY] - E[X]E[Y]}
\]

\subsection{Properties}
\begin{itemize}
    \item \textbf{Symmetry:} $\text{Cov}(X,Y)=\text{Cov}(Y,X)$
    \item \textbf{Variance case:} $\text{Cov}(X,X)=\text{Var}(X)$
    \item \textbf{Linearity:} $\text{Cov}(aX+b, cY+d)=ac \cdot \text{Cov}(X,Y)$
    \item \textbf{Independence:} $X \perp Y \Rightarrow \text{Cov}(X,Y)=0$
\end{itemize}
\textbf{Warning:} Zero covariance $\neq$ independence (in the general case).

\subsection{Example 1}
\textbf{Given:} $Y = -2X + 3$

\textbf{Step 1: Formula}
\[
\text{Cov}(X,Y)=E[XY] - E[X]E[Y]
\]

\textbf{Step 2: Compute } $XY$
\[
XY = X(-2X+3) = -2X^2 + 3X
\]

\textbf{Step 3: Compute } $E[Y]$
\[
E[Y] = E[-2X+3] = -2E[X] + 3
\]

\textbf{Step 4: Substitute}
\begin{align*}
\text{Cov}(X,Y) &= E[-2X^2+3X] - E[X](-2E[X] + 3) \\
&= -2E[X^2] + 3E[X] + 2(E[X])^2 - 3E[X] \\
&= -2E[X^2] + 2(E[X])^2
\end{align*}

\textbf{Step 5: Use variance}
\[
\text{Var}(X)=E[X^2]-(E[X])^2
\]
\[
\Rightarrow \text{Cov}(X,Y) = -2\text{Var}(X)
\]

\section{Pearson's Correlation Coefficient}

\subsection{Definition}
\[
\rho_{XY} = \frac{\text{Cov}(X,Y)}{\sigma_X \sigma_Y}
\]

\subsection{Interpretation}
\begin{itemize}
    \item \textbf{Measures:} Direction and Strength
    \item \textbf{Properties:} Unit-free, Normalized covariance
    \item \textbf{Range:} $-1 \le \rho_{XY} \le 1$
\end{itemize}

\subsection{Meaning}
\begin{itemize}
    \item $\rho > 0$: positive linear trend
    \item $\rho < 0$: negative linear trend
    \item $\rho = 0$: no linear correlation
\end{itemize}

\section{Joint Gaussian Random Variables}

\subsection{Definition}
A pair $(X,Y)$ is jointly Gaussian if:
\begin{itemize}
    \item Joint PDF has bivariate Gaussian form
    \item Marginals are Gaussian
\end{itemize}

\subsection{Algebraic Form}
\[
f_{XY}(x,y) =
\frac{1}{2\pi \sigma_X \sigma_Y \sqrt{1-\rho^2}}
\exp\left(
-\frac{1}{2(1-\rho^2)}
\left[
\left(\frac{x-\mu_X}{\sigma_X}\right)^2
+
\left(\frac{y-\mu_Y}{\sigma_Y}\right)^2
- 2\rho \left(\frac{x-\mu_X}{\sigma_X}\right)\left(\frac{y-\mu_Y}{\sigma_Y}\right)
\right]
\right)
\]

\subsection{Parameters}
\begin{itemize}
    \item $\mu_X, \mu_Y$: means
    \item $\sigma_X, \sigma_Y$: standard deviations
    \item $\rho$: correlation coefficient
\end{itemize}

\subsection{Interpretation (from lecture example)}
Example (page 23):
\begin{itemize}
    \item $X$: CO\textsubscript{2} concentration
    \item $Y$: temperature
\end{itemize}
If $\rho > 0 \Rightarrow$ Higher CO\textsubscript{2} $\rightarrow$ higher temperature.

\subsection{Applications}
\begin{itemize}
    \item Signal processing
    \item Communication systems
    \item Modeling noisy measurements
\end{itemize}

\vspace{1em}
\hrule
\vspace{1em}
\noindent \textbf{Final Concept Flow (Important for Exam)} \\
Joint PMF $\rightarrow$ Marginal PMF (via summation) $\rightarrow$ Product moment ($E[XY]$) $\rightarrow$ raw correlation $\rightarrow$ Centering $\rightarrow$ Covariance $\rightarrow$ Normalize $\rightarrow$ Pearson coefficient $\rightarrow$ Continuous extension $\rightarrow$ Joint Gaussian

\end{document}